\begin{abstract}
\noindent Derivatives on the National Stock Exchange of India settle against the
volume-weighted average price of the last thirty minutes of trading. We ask whether this window
is used to move the settlement price, using order-level data for every session of 2022,
calendar controls, a placebo group of securities without derivatives and a placebo index. The
answer depends on how the derivative is settled. Stock derivatives were settled by delivery,
which leaves a holder indifferent to the settlement price: the expiring stock future is priced
on the cash price until the end of the window, and the stock settlement shows no displacement
specific to securities with derivatives. Index derivatives were settled in cash, and on the
\IdxExpiries{} index expiries of the year the Nifty 50 falls over the settlement window by
\IdxSpreadDriftPoints{} index points relative to the Nifty Next 50, which has no derivatives
(randomization $p = \IdxSpreadDriftExpiryRip{}$). The fall occurs within the window and in the
index's constituents rather than in other stocks, and is followed the next morning by a
recovery of similar size, which is imprecisely estimated. Expiry also draws trading into the
window of securities with derivatives and widens their spreads, while order imbalance becomes
more persistent in all securities on monthly expiries.

\vspace{0.6em}
\noindent \textbf{Keywords:} expiration-day effects; settlement price; closing price
manipulation; order imbalance; physical settlement. \\
\noindent \textbf{JEL classification:} G13, G14, G18.
\end{abstract}
